\chapter{Tableaux de calibration}
\label{ann:calibration}
% BUDGET : 2,5 pages. Version resserree du 2026-08-31 : les six corrections de comparaison a
% l'article, le balayage plantain et ce que le solveur n'explique pas sont passes en tableaux.
% CORRIGE le 2026-08-31 : le plantain au seuil LE PLUS SERRE vaut 179 ha, pas 247 (qui est la
% valeur au seuil retenu de 6 440 t). Source : journal-vigilance 407-421, item CAL-76.
% Tous les nombres sont des macros lues dans les recap.json (build_chiffres.py), sauf le
% tableau du balayage plantain, saisi depuis CAL-76.

Le chapitre~\ref{ch:calibrer} affirme trois choses que seule cette annexe établit. L'écart
résiduel se concentre sur les petits postes. Le plafond de plantain et le plancher de prairie
sont des corrections et non des ajustements. Le solveur est hors de cause. Une section leur
répond dans cet ordre.

\section{\glsentryshort{pad} par culture, aux trois barreaux}

\begin{table}[h]
  \centering
  \caption{Surfaces observées et simulées par groupe \gls{rpg} (ha) et \gls{pad}
  correspondant (\%), aux trois barreaux de calibration. Le symbole $\bullet$ marque les
  cultures sous le seuil de \SI{\mesArtPad}{\percent} de \textcite{chopin2015}.
  \SI{\obsSurfaceIrreproductible}{\ha} du territoire ne sont reproductibles par aucune
  variante fine, ce qui impose un plancher d'écart de \num{\obsPlancherPad} point avant toute
  optimisation.}
  \label{tab:ann-pad-cultures}
  \scriptsize
  \begin{tabular}{@{}lrrrrrrr@{}}
    \toprule
    & & \multicolumn{2}{c}{\textbf{Parité \glsentryshort{gams}}}
      & \multicolumn{2}{c}{\textbf{+ plafond plantain}}
      & \multicolumn{2}{c}{\textbf{+ plancher prairie}} \\
    \cmidrule(lr){3-4}\cmidrule(lr){5-6}\cmidrule(lr){7-8}
    \textbf{Groupe} & \textbf{observé} & sim. & \gls{pad} & sim. & \gls{pad}
      & sim. & \gls{pad} \\
    \midrule
    Canne à sucre (CS) & \num{\obsHaCS} & \num{\simPariteCS} & \num{\padPariteCS}
      & \num{\simBcSeulCS} & \num{\padBcSeulCS} & \num{\simRetenuCS}
      & \num{\padRetenuCS}~$\bullet$ \\
    Prairie (PN) & \num{\obsHaPN} & \num{\simParitePN} & \num{\padParitePN}
      & \num{\simBcSeulPN} & \num{\padBcSeulPN} & \num{\simRetenuPN}
      & \num{\padRetenuPN}~$\bullet$ \\
    Banane export (BA) & \num{\obsHaBA} & \num{\simPariteBA} & \num{\padPariteBA}
      & \num{\simBcSeulBA} & \num{\padBcSeulBA} & \num{\simRetenuBA}
      & \num{\padRetenuBA}~$\bullet$ \\
    Maraîchage (MA) & \num{\obsHaMA} & \num{\simPariteMA} & \num{\padPariteMA}
      & \num{\simBcSeulMA} & \num{\padBcSeulMA} & \num{\simRetenuMA}
      & \num{\padRetenuMA}~$\bullet$ \\
    Jachère (JA) & \num{\obsHaJA} & \num{\simPariteJA} & \num{\padPariteJA}
      & \num{\simBcSeulJA} & \num{\padBcSeulJA} & \num{\simRetenuJA} & \num{\padRetenuJA} \\
    Plantain (BC) & \num{\obsHaBC} & \num{\simPariteBC} & \num{\padPariteBC}
      & \num{\simBcSeulBC} & \num{\padBcSeulBC} & \num{\simRetenuBC} & \num{\padRetenuBC} \\
    Igname (IG) & \num{\obsHaIG} & \num{\simPariteIG} & \num{\padPariteIG}
      & \num{\simBcSeulIG} & \num{\padBcSeulIG} & \num{\simRetenuIG} & \num{\padRetenuIG} \\
    Ananas (AN) & \num{\obsHaAN} & \num{\simPariteAN} & \num{\padPariteAN}
      & \num{\simBcSeulAN} & \num{\padBcSeulAN} & \num{\simRetenuAN} & \num{\padRetenuAN} \\
    Melon (ME) & \num{\obsHaME} & \num{\simPariteME} & \num{\padPariteME}
      & \num{\simBcSeulME} & \num{\padBcSeulME} & \num{\simRetenuME} & \num{\padRetenuME} \\
    Agrumes (AG) & \num{\obsHaAG} & \num{\simPariteAG} & \num{\padPariteAG}
      & \num{\simBcSeulAG} & \num{\padBcSeulAG} & \num{\simRetenuAG} & \num{\padRetenuAG} \\
    Vergers (VE) & \num{\obsHaVE} & \num{\simPariteVE} & \num{\padPariteVE}
      & \num{\simBcSeulVE} & \num{\padBcSeulVE} & \num{\simRetenuVE} & \num{\padRetenuVE} \\
    \midrule
    cultures sous le seuil & & & \num{\calibPariteCulturesOk}
      & & \num{\calibBcSeulCulturesOk} & & \num{\calibRetenuCulturesOk} \\
    \bottomrule
  \end{tabular}
\end{table}

Cette table établit trois choses. Les quatre cultures qui passent sous le seuil sont les
quatre plus grandes en surface, et les sept qui échouent totalisent moins d'un dixième du
territoire. Le plancher d'éligibilité joue ensuite dans les deux sens, puisqu'il innocente le
modèle sur les vergers, dont une partie n'est accessible à aucune variante fine, et l'accuse
partout ailleurs, où rien n'interdisait de faire mieux. Le succès sur la prairie, enfin, est
nul par construction puisque son plancher l'épingle, à l'inverse exact de la canne, qu'aucune
contrainte ne nomme.

\section{Les autres échelles, et la comparaison à l'article}

Le tableau~\ref{tab:ann-echelles} donne les trois autres échelles, à la parité et sous la
calibration retenue.

\begin{table}[h]
  \centering
  \caption{Les trois autres échelles d'évaluation. Le seuil sous-régional et le seuil par
  ferme valent \SI{\mesArtPadZone}{\percent} et non \SI{\mesArtPad}{\percent}.}
  \label{tab:ann-echelles}
  \scriptsize
  \setlength{\tabcolsep}{4pt}
  \begin{tabularx}{\textwidth}{@{}p{28mm}rrX@{}}
    \toprule
    \textbf{Échelle} & \textbf{parité} & \textbf{retenu} & \textbf{Lecture} \\
    \midrule
    Sous-régions sous le seuil & & \num{\calibRetenuRegionsOk}/\num{\calibRetenuRegionsTotal} &
      les reçues tiennent entre \SI{\calibRetenuRegionsPadMin}{\percent} et
      \SI{\calibRetenuRegionsPadOkMax}{\percent}. Les deux échecs sont le Sud-Est et le
      Sud-Ouest de Basse-Terre, à \SI{\calibRetenuRegionsPadHorsMin}{\percent} et
      \SI{\calibRetenuRegionsPadMax}{\percent} \\
    Fermes sous le seuil (\%) & &
      \num{\calibRetenuFermesOkPct} &
      \num{\calibRetenuFermesOk} fermes sur \num{\calibRetenuFermesTotal} \\
    \gls{pad} par ferme, médiane (\%) & \num{\calibPariteFermesPadMediane} &
      \num{\calibRetenuFermesPadMediane} &
      plus d'une ferme sur deux est reproduite exactement \\
    \gls{pad} par ferme, moyenne (\%) & \num{\calibPariteFermesPadMoyenne} &
      \num{\calibRetenuFermesPadMoyenne} &
      portée par une minorité de fermes très mal simulées. La lire comme un écart uniforme
      sous-estimerait largement ce que le modèle reproduit \\
    Types reproduits (\%) & \num{\calibPariteTypes} & \num{\calibRetenuTypes} &
      matrice de confusion des huit types d'exploitation \\
    Parcelles justes (\%) & & \num{\calibRetenuParcelles} &
      pour \SI{\calibRetenuSurface}{\percent} de la surface. Le modèle se trompe donc
      davantage sur les petites parcelles \\
    \bottomrule
  \end{tabularx}
\end{table}

Deux lectures s'ajoutent à cette table. L'article échoue sur la même sous-région que nous, le
Sud-Est de Basse-Terre, sa seule à descendre à \SI{56}{\percent} de surface correctement
simulée, et il en donne la cause. La rotation banane et canne forme un système technique
unique réparti sur plusieurs parcelles, que le modèle traite comme des choix indépendants.
Marie-Galante, elle, est passée du pire au meilleur rang, sa prairie tombant de
\SI{1357}{\ha} à \SI{228}{\ha} avant le plancher, soit \SI{83}{\percent} de perte sur une île
entière. L'écart n'est pas diffus, il a une adresse.

Six corrections sont enfin nécessaires avant de rapprocher ces chiffres de ceux de l'article,
et le tableau~\ref{tab:ann-corrections} les énonce.

\begin{table}[h]
  \centering
  \caption{Les six corrections nécessaires pour comparer honnêtement ces chiffres à
  \textcite{chopin2015}. Elles doivent être énoncées avant toute conclusion.}
  \label{tab:ann-corrections}
  \scriptsize
  \setlength{\tabcolsep}{4pt}
  \begin{tabularx}{\textwidth}{@{}p{32mm}X@{}}
    \toprule
    \textbf{Point} & \textbf{Ce qu'il impose} \\
    \midrule
    Le run de parité n'est pas le modèle validé par l'article &
      l'article valide un modèle qui contient déjà le plafond de plantain, son équation~6.
      Notre parité le désactive pour mesurer ce qu'il apporte, et c'est la cause dominante de
      l'écart, le plantain y montant à \SI{\simPariteBC}{\ha} contre \SI{\obsHaBC}{\ha} \\
    L'article ne publie aucun \gls{pad} agrégé &
      la seule métrique comparable est le nombre de cultures sous le seuil, celle sur laquelle
      notre modèle est le plus faible. Ses dix usages fusionnent vergers et agrumes, les deux
      postes que nous ratons le plus \\
    Ses seuils sont deux et non un &
      \SI{\mesArtPad}{\percent} au territoire, \SI{\mesArtPadZone}{\percent} en sous-région et
      à la ferme. Comparer nos chiffres de ferme au premier serait gratuitement défavorable \\
    Les tables parcellaires n'ont pas le même dénominateur &
      la sienne compte les parcelles non cultivées, la nôtre non. Quatre points d'écart de
      définition \\
    Son année de base est 2010 &
      \num{\mesArtParcellesN} parcelles et \num{\mesArtFermes} exploitations contre
      \num{\obsParcelles} et \num{\obsExploitations} ici, soit \SI{13}{\percent}
      d'exploitations en moins à nombre de parcelles quasi identique. Les coefficients
      d'aversion de sa Table~2 étant calibrés sur ces fermes, les appliquer à 2017 est un
      transfert hors échantillon \\
    L'article ne publie aucune table par exploitation &
      notre quatrième échelle n'a pas de référence et ne se lit qu'en absolu \\
    \bottomrule
  \end{tabularx}
\end{table}

\section{Les balayages de seuil}

Sans ces deux balayages, \enquote{correction et non ajustement} reste une affirmation. Le
tableau~\ref{tab:ann-balayage-plantain} donne le premier en entier.

\begin{table}[h]
  \centering
  \caption{Balayage du plafond de plantain, sept seuils. Le \gls{pad} y forme un palier plat
  de \SI{4650}{\tonne} à \SI{9240}{\tonne}, soit un facteur deux sur le seuil pour un
  demi-point, à peine plus que le bruit du solveur. Ces valeurs datent du balayage du 27
  juillet, avant l'activation du plancher de prairie. Leur niveau n'est donc pas comparable au
  tableau~\ref{tab:ann-pad-cultures}, seule leur variation le long du balayage est en cause.}
  \label{tab:ann-balayage-plantain}
  \small
  \setlength{\tabcolsep}{4pt}
  \begin{tabular}{@{}lrrrrrrr@{}}
    \toprule
    \textbf{Seuil (t)} & \num{4650} & \num{6440} & \num{9240} & \num{15000} & \num{25000}
      & \num{40000} & désactivé \\
    \midrule
    \gls{pad} (\%) & \num{31.7} & \num{31.7} & \num{31.2} & \num{34.4} & \num{36.4}
      & \num{38.1} & \num{48.4} \\
    Types (\%) & \num{66.8} & \num{67.0} & \num{67.1} & \num{65.8} & \num{65.6} & \num{65.4}
      & \num{64.0} \\
    Surface (\%) & \num{69.0} & \num{69.0} & \num{69.3} & \num{68.6} & \num{68.2} & \num{67.7}
      & \num{64.3} \\
    Plantain (ha) & \num{179} & \num{247} & \num{355} & \num{577} & \num{960} & \num{1527}
      & \num{2875} \\
    \bottomrule
  \end{tabular}
\end{table}

Deux observations écartent le soupçon de circularité. L'optimum du palier se trouve à
l'extrémité haute, à deux fois et demie le débouché observé, et non au seuil le plus serré,
alors qu'un paramètre servant de variable d'ajustement s'améliorerait de façon monotone en se
rapprochant de l'observé. Et le plantain simulé reste au-dessus de l'observé
% CORRIGE 2026-08-31 : la version precedente ecrivait 247 ha ici. 247 ha est la valeur au
% seuil RETENU (6 440 t) ; au seuil le plus serre (4 650 t) c'est 179 ha. Source CAL-76.
(\SI{\obsHaBC}{\ha}) à tous les seuils, le minimum étant \SI{179}{\ha} au plus serré. La
contrainte ne pousse donc pas la culture vers l'observé. Le seuil retenu est
\SI{6440}{\tonne}, qui n'est pas celui qui note le mieux mais le seul des trois candidats qui
soit sourçable.

Sur la prairie, le seuil a été balayé de sa valeur retenue jusqu'à la statistique exogène. Les
trois métriques que le plancher ne touche pas restent au-dessus des seuils publiés sur toute
la plage, avec au plancher le plus haut \SI{\mesPrairieHautTypes}{\percent} de types,
\SI{\mesPrairieHautParcelles}{\percent} de parcelles et \SI{\mesPrairieHautSurface}{\percent}
de surface. Le \gls{pad} hors prairie empire quand on remonte le seuil, de
\SI{\mesPadHorsPrairieBas}{\percent} à \SI{\mesPadHorsPrairieHaut}{\percent}, parce que le
plancher prend alors sa terre à la canne, la culture la mieux restituée. Un seuil calé pour
flatter le résultat ferait exactement l'inverse.

Le meilleur argument du chapitre~\ref{ch:calibrer} est cependant extérieur au modèle. Le
cheptel de \num{\mesBovins} bovins recensés en 2017 vaut \num{\mesUgb} \gls{ugb}. Rapportés
aux \SI{\mesPrairieAgreste}{\ha} de la statistique agricole, cela fait
\num{\mesUgbHaAgreste} \gls{ugb} par hectare, conforme aux recensements. Rapportés aux
\SI{2980}{\ha} que le modèle sans plancher laisse à la prairie, \num{\mesUgbHaSansPlancher},
agronomiquement impossible. La prairie de notre jeu parcellaire couvre
\SI{\mesPrairieRpg}{\ha}, soit \SI{\mesPrairieExogene}{\ha} une fois la statistique ramenée à
notre couverture de \SI{\mesCouvertureSau}{\percent}. Sans plancher, \SI{3489}{\ha} de
parcelles observées en prairie ne sont pas reconduites, dont \SI{62}{\percent} d'économie
légitime, \SI{36}{\percent} de couplage avec le plafond de main d'\oe uvre et
\SI{2}{\percent} de résidu. \SI{30}{\percent} du déficit tient sur un écart de
\SI{0.94}{\percent} entre les marges de deux cultures, et la prairie pesait
\SI{42}{\percent} de la déviation totale avant son plancher.

\section{Ce que le solveur n'explique pas}

Cinq suspects ont été instruits et écartés, et deux verdicts que j'avais tenus ont été
démentis. Le tableau~\ref{tab:ann-solveur} donne pour chacun la mesure qui a tranché.

\begin{table}[h]
  \centering
  \caption{Ce qui a été mis hors de cause, et par quelle mesure. Les deux dernières lignes
  sont des verdicts que j'ai tenus puis démentis, et leur date figure en
  annexe~\ref{ann:chronologie}.}
  \label{tab:ann-solveur}
  \scriptsize
  \setlength{\tabcolsep}{4pt}
  \begin{tabularx}{\textwidth}{@{}p{34mm}X@{}}
    \toprule
    \textbf{Suspect} & \textbf{Ce qui l'écarte} \\
    \midrule
    L'arbitraire du solveur & le même problème résolu avec trois graines de branchement rend
      \SI{\mesGraineA}{\percent}, \SI{\mesGraineB}{\percent} et \SI{\mesGraineC}{\percent} de
      \gls{pad} territorial. L'arbitraire vaut donc \num{\mesBruitPad} point, et c'est ce
      chiffre qui fixe à \SI{5}{\percent} la tolérance du garde-fou des runs de référence \\
    La main d'\oe uvre & le plan réellement observé demande \SI{5.68}{\mega\hour} pour un
      plafond de \SI{6.25}{\mega\hour}. Le modèle n'a pas d'excuse de ce côté \\
    Recalibrer l'aversion au risque & descend le \gls{pad} à \SI{39.9}{\percent} avec des
      coefficients économiquement absurdes. Surajustement \\
    Caler des plafonds de marché sur l'observé & descend à \SI{33.6}{\percent} et c'est
      circulaire par construction \\
    Incorporer l'élevage & c'est la bonne réponse, et c'est pourquoi elle est écartée du stage
      plutôt que bâclée. Elle demande une variable d'état que le modèle mono-période ne porte
      pas \\
    \enquote{Le plancher de prairie rend le \gls{milp} intraitable} & verdict daté d'une
      version du modèle à \num{900000} variables. Il n'a pas survécu au masque
      d'éligibilité \\
    \enquote{Le plancher de prairie est circulaire} & mon propre verdict, tenu jusqu'au jour
      où la donnée exogène de charge animale a été trouvée \\
    \bottomrule
  \end{tabularx}
\end{table}

Une dernière réserve porte sur les indicateurs plutôt que sur les surfaces. L'assolement 2017
n'étant connu qu'au niveau des douze groupes agrégés, chaque indicateur observé est une
fourchette et non une valeur. Plusieurs indicateurs simulés, la subvention, le revenu, l'azote
et le travail, tombent à l'intérieur de cette fourchette, et le point central en sort parfois
par le haut. L'écart au central ne démontre donc rien, et c'est une réserve à opposer à toute
comparaison indicateur par indicateur.
